% ==================================================
% type_system.tex
% TYPE SYSTEM LAYER (safe, no side effects)
% Only defines reusable macros; does NOT patch Sphinx styles directly.
% ==================================================

% --- helpers ---
\providecommand{\HBifupper}[2]{%
  \ifnum#1=1\relax\MakeUppercase{#2}\else#2\fi
}

% ==================================================
% TYPE PRESETS (macros only)
% ==================================================

% Base document text.
\providecommand{\HBTypeBaseStart}{%
  \HBFontRegular
  \color{\csname HBtype_body_color\endcsname}%
  \fontsize{\csname HBtype_body_font_size\endcsname}{\csname HBtype_body_font_leading\endcsname}\selectfont
}

% H1 title in dark pill
\providecommand{\HBTypeHOne}[1]{%
  {\color{white}\HBFontHeavy
   \fontsize{\csname HBtype_h1_font_size\endcsname}{\csname HBtype_h1_font_leading\endcsname}\selectfont
   \HBifupper{\csname HBtype_h1_force_upper\endcsname}{#1}}%
}

% Structural title level 1. Keep as an alias so existing H1 params remain stable.
\providecommand{\HBTypeTitleLevelOne}[1]{%
  \HBTypeHOne{#1}%
}

% Structural title level 2. Used by RST subsections, overview panels,
% spec-section titles, and subbar-like titles.
\expandafter\providecommand\csname HBtype_title_l2_color\endcsname{TextGray}
\expandafter\providecommand\csname HBtype_title_l2_font_size\endcsname{8.6pt}
\expandafter\providecommand\csname HBtype_title_l2_font_leading\endcsname{9.4pt}
\expandafter\providecommand\csname HBtype_title_l2_force_upper\endcsname{1}
\providecommand{\HBTypeTitleLevelTwo}[1]{%
  {\HBFontHeavy
   \color{\csname HBtype_title_l2_color\endcsname}%
   \fontsize{\csname HBtype_title_l2_font_size\endcsname}
            {\csname HBtype_title_l2_font_leading\endcsname}\selectfont
   \RaggedRight
   \HBifupper{\csname HBtype_title_l2_force_upper\endcsname}{#1}}%
}

% Structural title level 3. Used by RST subsubsections and lower local
% headings.
\expandafter\providecommand\csname HBtype_title_l3_color\endcsname{TextGray}
\expandafter\providecommand\csname HBtype_title_l3_font_size\endcsname{7.0pt}
\expandafter\providecommand\csname HBtype_title_l3_font_leading\endcsname{8.0pt}
\expandafter\providecommand\csname HBtype_title_l3_force_upper\endcsname{0}
\providecommand{\HBTypeTitleLevelThree}[1]{%
  {\HBFontSemiBold
   \color{\csname HBtype_title_l3_color\endcsname}%
   \fontsize{\csname HBtype_title_l3_font_size\endcsname}
            {\csname HBtype_title_l3_font_leading\endcsname}\selectfont
   \RaggedRight
   \HBifupper{\csname HBtype_title_l3_force_upper\endcsname}{#1}}%
}

% Subbar white text on dark background
\providecommand{\HBTypeSubbar}[1]{%
  {\color{white}\HBFontSemiBold
   \fontsize{\csname HBtype_subbar_font_size\endcsname}{\csname HBtype_subbar_font_leading\endcsname}\selectfont
   \HBifupper{\csname HBtype_subbar_force_upper\endcsname}{#1}}%
}

% Body/lead
\providecommand{\HBTypeBody}[1]{%
  {\HBFontMedium
   \color{\csname HBtype_body_color\endcsname}%
   \fontsize{\csname HBtype_body_font_size\endcsname}{\csname HBtype_body_font_leading\endcsname}\selectfont
   \HBifupper{\csname HBtype_body_force_upper\endcsname}{#1}}%
}

% List (used in two-column env)
\providecommand{\HBTypeListStart}{%
  \HBFontRegular
  \color{\csname HBtype_list_color\endcsname}%
  \fontsize{\csname HBtype_list_font_size\endcsname}{\csname HBtype_list_font_leading\endcsname}\selectfont
}

% Warning right text (font only; alignment handled elsewhere)
\providecommand{\HBTypeWarningTextStart}{%
  \HBFontSemiBold
  \color{\csname HBtype_warning_text_color\endcsname}%
  \fontsize{\csname HBtype_warning_text_font_size\endcsname}{\csname HBtype_warning_text_font_leading\endcsname}\selectfont
}

% Rubric header (font only; sectioning handled in components_base.tex)
\providecommand{\HBTypeRubricStart}{%
  \HBFontHeavy
  \fontsize{\csname HBtype_rubric_font_size\endcsname}{\csname HBtype_rubric_font_leading\endcsname}\selectfont
  \color{\csname HBtype_rubric_color\endcsname}%
}

% ==================================================
% FOOTER TYPE
% ==================================================

\newcommand{\HBTypeFooter}[1]{%
  {\color{\csname HBpage_footer_color\endcsname}%
   \fontsize{\csname HBpage_footer_font_size\endcsname}
            {\csname HBpage_footer_font_leading\endcsname}%
   \selectfont
   #1}%
}

% ==================================================
% PAGE NUMBER TYPE
% ==================================================

\newcommand{\HBTypePageNumber}{%
  \ifnum\csname HBpage_footer_leading_zero\endcsname=1\relax
    \ifnum\value{page}<10\relax 0\fi
  \fi
  \arabic{page}%
}
